% =============================================================================
%  Capítulo 1 — Antes de Começar (Parte I, Visão Geral)
%  Capítulo único da parte: público, requisitos de acesso, como ler o guia e
%  o fluxo de uso em BPMN, com a legenda e o mapa passo -> tela -> página.
% =============================================================================
\chapter{Antes de Começar}

\section{Público e escopo}

Este guia se destina ao usuário com perfil \textbf{Gestor}, que mantém a base de
scrap atualizada e acompanha o IF Cost da Linha Final. A leitura completa leva cerca
de quinze minutos.

\begin{itemize}
    \item \textbf{Coberto:} acesso, envio do relatório do GERP, conferência da
          execução, Dashboard, Base de Scrap, análise de uma ocorrência, aplicação de
          modelo em lote e Modo TV.
    \item \textbf{Não coberto:} relatórios de análise, alertas e regras de
          notificação, configurações (classificações de material, tipos de scrap,
          programa do ano) e administração de usuários.
\end{itemize}

\section{Requisitos de acesso}

\begingroup
\guiazebra
\small
\noindent\begin{longtable}{@{}L{4.0cm}L{11.5cm}@{}}
    \toprule
    \guiacabtabela
    \thead{Item} & \thead{Requisito} \\
    \midrule
    \endhead
    \bottomrule
    \endlastfoot
    Navegador           & Navegador atualizado. As telas deste guia foram capturadas no
                          Google Chrome \\
    Endereço do sistema & Informado pela equipe na liberação do acesso \\
    Credenciais         & Usuário e senha do Hanaro, cadastrados pelo administrador \\
    Perfil              & Gestor \\
    Arquivo de entrada  & Relatório \emph{Other Account Transaction Text} exportado do
                          GERP, em texto separado por tabulação (TSV), com codificação
                          UTF-8 ou CP1252 \\
    Modo TV             & TV ou monitor da linha com navegador, para a exibição em tela
                          cheia \\
\end{longtable}
\endgroup

\begin{remark}
O usuário só entra no Hanaro depois de cadastrado e ativado pelo administrador do
sistema. Sem esse cadastro, a tela de acesso recusa as credenciais.
\end{remark}

\section{Como ler este guia}

O \autoref{cap:passos} descreve o uso em \textbf{passos}. Cada passo ocupa uma linha
da página, dividida em duas colunas:

\begin{itemize}
    \item \textbf{à esquerda, a tela} do sistema, exatamente como o usuário a vê;
    \item \textbf{à direita, o texto} do passo: o que fazer, em que ordem e o que o
          sistema responde.
\end{itemize}

Os números em vermelho sobre a tela, como \marca{1} e \marca{2}, indicam os elementos
citados no texto ao lado, na mesma ordem. O número do passo, no selo
\colorbox{guiavermelho}{\color{white}\guiahead\bfseries\fontsize{7.5}{7.5}\selectfont PASSO 1},
é o mesmo que aparece no diagrama BPMN da \autoref{sec:fluxo}.

Botões e opções da interface aparecem em destaque, como \botao{Entrar}. Um caminho de
menu é escrito como \caminho{Menu\seta Execuções}.

As caixas ao longo do texto têm três tipos:

\begin{remark}
\textbf{Atenção.} Condição que impede o passo de ser concluído ou que exige
correção.
\end{remark}

\begin{remark}[dica]
\textbf{Dica.} Atalho ou forma mais rápida de concluir o passo.
\end{remark}

\begin{remark}[info]
\textbf{Informação.} Comportamento do sistema que convém conhecer, sem ação
necessária.
\end{remark}

% ---------------------------------------------------------------------------
%  Fluxo de uso em BPMN
%
%  O diagrama usa os estilos `bpmn ...' do axacademy-guia.sty. As raias são
%  COLUNAS e o fluxo desce: assim o diagrama cabe em pé, nos 16 cm da largura
%  do texto. Coordenadas em cm, com y = 0 no topo das raias.
%
%  A seção abre página nova: o diagrama vai em `figure[H]', que não flutua, e
%  precisa da página inteira. Com oito passos, as tarefas têm uma linha só e
%  as linhas do diagrama ficam a 1,15 cm uma da outra; com mais passos,
%  divida o fluxo em dois diagramas.
% ---------------------------------------------------------------------------
\clearpage
\section{Fluxo de uso}\label{sec:fluxo}

O fluxo básico vai do acesso ao Modo TV. A raia \textbf{Gestor} reúne o que o
usuário faz; a raia \textbf{Hanaro}, o que o sistema faz em resposta. As tarefas com
borda vermelha são os passos deste guia: o número no canto é o do passo e leva à
página em que ele é descrito.

\begin{figure}[H]
\centering
\begin{tikzpicture}[x=1cm,y=1cm,
        bpmn tarefa/.append style={text width=4.4cm,minimum height=0.72cm}]
    % A caixa de medida é a das raias, exata: sem ela, a espessura do contorno
    % passa 0,7 pt da largura do texto.
    \useasboundingbox (0,0) rectangle (16,-13.35);
    % ---- raias ----
    \bpmnraia{0}{8}{13.35}{Gestor}
    \bpmnraia{8}{8}{13.35}{Hanaro}

    % ---- gestor ----
    \node[bpmn inicio] (ini) at (4,-1.05) {};
    \node[bpmn rotulo,anchor=west] at ([xshift=4pt]ini.east)
        {Relatório do GERP\\disponível};
    \bpmnpasso{p1}{4,-2.15}{1}{Entrar com usuário e senha}
    \bpmnpasso{p2}{4,-4.45}{2}{Enviar o relatório do GERP}
    \bpmnpasso{p3}{4,-6.75}{3}{Conferir a execução}
    \bpmnpasso{p4}{4,-7.90}{4}{Ler o Dashboard}
    \bpmnpasso{p5}{4,-9.05}{5}{Revisar a Base de Scrap}
    \bpmnpasso{p6}{4,-10.20}{6}{Consultar a análise do caso}
    \bpmnpasso{p7}{4,-11.35}{7}{Aplicar o modelo em lote}
    \bpmnpasso{p8}{4,-12.50}{8}{Exibir o Modo TV}

    % ---- sistema ----
    \node[bpmn tarefa] (s1) at (12,-2.15) {Validar as credenciais};
    \bpmngateway{g1}{12,-3.30}
    \node[bpmn rotulo,anchor=west] at ([xshift=3pt]g1.east) {Credenciais\\válidas?};
    \node[bpmn tarefa] (s2) at (12,-4.45) {Abrir o Dashboard};
    \node[bpmn tarefa] (s3) at (12,-5.60) {Processar o arquivo enviado};
    \node[bpmn tarefa] (s4) at (12,-6.75) {Publicar o snapshot da base};
    \node[bpmn tarefa] (s5) at (11.6,-12.50) {Girar as telas do Modo TV};
    \node[bpmn fim] (fim) at (14.95,-12.50) {};
    \node[bpmn rotulo,anchor=south,align=center] at ([yshift=3pt]fim.north)
        {Indicadores\\na fábrica};

    % ---- fluxos ----
    \draw[bpmn fluxo] (ini) -- (p1);
    \draw[bpmn fluxo] (p1) -- (s1);
    \draw[bpmn fluxo] (s1) -- (g1);
    \draw[bpmn fluxo] (g1) -- node[bpmn rotulo,anchor=west,xshift=2pt] {Sim} (s2);
    \draw[bpmn fluxo] (g1.west) -- node[bpmn rotulo,anchor=south,pos=0.25] {Não}
        (4,-3.30) -- (p1.south);
    \draw[bpmn fluxo] (s2) -- (p2);
    \draw[bpmn fluxo] (p2.south) |- (s3.west);
    \draw[bpmn fluxo] (s3) -- (s4);
    \draw[bpmn fluxo] (s4) -- (p3);
    \draw[bpmn fluxo] (p3) -- (p4);
    \draw[bpmn fluxo] (p4) -- (p5);
    \draw[bpmn fluxo] (p5) -- (p6);
    \draw[bpmn fluxo] (p6) -- (p7);
    \draw[bpmn fluxo] (p7) -- (p8);
    \draw[bpmn fluxo] (p8) -- (s5);
    \draw[bpmn fluxo] (s5) -- (fim);
\end{tikzpicture}
\caption{Fluxo básico de uso do Hanaro pelo Gestor, em BPMN. Os números indicam os
passos do \autoref{cap:passos}.}
\label{fig:bpmn}
\end{figure}

\section{Leitura do diagrama}

\begingroup
\small
\renewcommand{\arraystretch}{1.6}
\noindent\begin{tabularx}{\textwidth}{@{}M{2.4cm}X@{}}
    \toprule
    \guiacabtabela
    \thead{Símbolo} & \thead{Significado} \\
    \midrule
    \tikz[baseline=-0.6ex]\node[bpmn inicio,minimum size=0.42cm]{};
        & Início do fluxo: o evento que leva o usuário ao sistema. \\
    \tikz[baseline=-0.6ex]\node[bpmn tarefa,text width=1.2cm,minimum height=0.5cm,
            font=\fontsize{7}{8}\selectfont]{Tarefa};
        & Tarefa executada pelo usuário ou pelo sistema, conforme a raia em que está. \\
    \tikz[baseline=-0.6ex]{\node[bpmn tarefa,draw=guiavermelho,line width=1.1pt,
            text width=1.2cm,minimum height=0.5cm,font=\fontsize{7}{8}\selectfont] (t) {Passo};
            \node[guia marcador,minimum size=11pt,font=\guiahead\bfseries\fontsize{7}{7}\selectfont]
            at (t.north west) {3};}
        & Tarefa descrita neste guia. O número é o do passo; no PDF, ele leva à página
          correspondente. \\
    \tikz[baseline=-0.6ex]{\bpmngateway{g}{0,0}}
        & Decisão. O fluxo segue por um único caminho, conforme a condição escrita ao
          lado. \\
    \tikz[baseline=-0.6ex]\draw[bpmn fluxo] (0,0) -- (1.2,0);
        & Sequência: a próxima tarefa só começa quando a anterior termina. \\
    \tikz[baseline=-0.6ex]\node[bpmn fim,minimum size=0.42cm]{};
        & Fim do fluxo. \\
    \bottomrule
\end{tabularx}
\endgroup

\section{Mapa dos passos}

\begingroup
\guiazebra
\small
\noindent\begin{longtable}{@{}P{1.3cm}L{3.2cm}L{8.4cm}P{1.4cm}@{}}
    \toprule
    \guiacabtabela
    \thead{Passo} & \thead{Tela} & \thead{Ação principal} & \thead{Página} \\
    \midrule
    \endhead
    \bottomrule
    \endlastfoot
    \marca{1} & Acesso           & Informar usuário e senha e clicar em \botao{Entrar}       & \pagina{passo:1} \\
    \marca{2} & Execuções        & Selecionar o arquivo do GERP e clicar em
                                   \botao{Processar relatório}                               & \pagina{passo:2} \\
    \marca{3} & Execuções        & Conferir o status e o snapshot da execução                & \pagina{passo:3} \\
    \marca{4} & Dashboard        & Ler o IF Cost acumulado, o target e os rankings           & \pagina{passo:4} \\
    \marca{5} & Base de Scrap    & Filtrar e revisar as ocorrências de scrap                 & \pagina{passo:5} \\
    \marca{6} & Análise de scrap & Ler o custo, a contenção e as evidências da ocorrência    & \pagina{passo:6} \\
    \marca{7} & Base de Scrap    & Marcar as ocorrências e aplicar o modelo de análise       & \pagina{passo:7} \\
    \marca{8} & Modo TV          & Abrir o Modo TV e controlar a rotação das telas           & \pagina{passo:8} \\
\end{longtable}
\endgroup

\fimdecapitulo
